% Created 2025-10-30 jue 17:19
% Intended LaTeX compiler: pdflatex
\documentclass[a4paper]{article}
\usepackage[utf8]{inputenc}
\usepackage[T1]{fontenc}
\usepackage{graphicx}
\usepackage{longtable}
\usepackage{wrapfig}
\usepackage{rotating}
\usepackage[normalem]{ulem}
\usepackage{amsmath}
\usepackage{amssymb}
\usepackage{capt-of}
\usepackage{hyperref}
\usepackage[spanish]{babel}
\usepackage[usenames,dvipsnames]{color} % Required for custom colors
\renewcommand{\ttdefault}{pcr} % MONOESPACIO CON NEGRIT
\usepackage{lastpage}
\usepackage{listings}
\usepackage{listingsutf8}
\renewcommand{\lstlistingname}{Listado}
\lstset{frame=single,inputencoding=utf8,basicstyle=\scriptsize\ttfamily,showstringspaces=false,numbers=none}
\definecolor{MyDarkGreen}{rgb}{0.0,0.4,0.0} % This is the color used for comments
\lstset{ breaklines=true, postbreak=\mbox{\textcolor{red}{$\hookrightarrow$}\space}, keywordstyle=\bfseries, keywordstyle=[1]\color{Blue}\bfseries,  keywordstyle=[2]\color{Purple}\bfseries,  keywordstyle=[3]\color{Blue}\underbar,   identifierstyle=,   commentstyle=\usefont{T1}{pcr}{m}{sl}\color{MyDarkGreen}\small,   stringstyle=\color{Purple},   showstringspaces=false,   tabsize=2,   morecomment=[l][\color{Blue}]{...} }
\lstset{literate=  {á}{{\'a}}1 {é}{{\'e}}1 {í}{{\'i}}1 {ó}{{\'o}}1 {ú}{{\'u}}1   {Á}{{\'A}}1 {É}{{\'E}}1 {Í}{{\'I}}1 {Ó}{{\'O}}1 {Ú}{{\'U}}1   {à}{{\`a}}1 {è}{{\`e}}1 {ì}{{\`i}}1 {ò}{{\`o}}1 {ù}{{\`u}}1   {À}{{\`A}}1 {È}{{\'E}}1 {Ì}{{\`I}}1 {Ò}{{\`O}}1 {Ù}{{\`U}}1   {ä}{{\"a}}1 {ë}{{\"e}}1 {ï}{{\"i}}1 {ö}{{\"o}}1 {ü}{{\"u}}1   {Ä}{{\"A}}1 {Ë}{{\"E}}1 {Ï}{{\"I}}1 {Ö}{{\"O}}1 {Ü}{{\"U}}1   {â}{{\^a}}1 {ê}{{\^e}}1 {î}{{\^i}}1 {ô}{{\^o}}1 {û}{{\^u}}1   {Â}{{\^A}}1 {Ê}{{\^E}}1 {Î}{{\^I}}1 {Ô}{{\^O}}1 {Û}{{\^U}}1   {œ}{{\oe}}1 {Œ}{{\OE}}1 {æ}{{\ae}}1 {Æ}{{\AE}}1 {ß}{{\ss}}1   {ű}{{\H{u}}}1 {Ű}{{\H{U}}}1 {ő}{{\H{o}}}1 {Ő}{{\H{O}}}1   {ç}{{\c c}}1 {Ç}{{\c C}}1 {ø}{{\o}}1 {å}{{\r a}}1 {Å}{{\r A}}1   {€}{{\euro}}1 {£}{{\pounds}}1 {«}{{\guillemotleft}}1   {»}{{\guillemotright}}1 {ñ}{{\~n}}1 {Ñ}{{\~N}}1 {¿}{{?`}}1 }
\usepackage{caption}
\usepackage{attachfile2}
\usepackage[margin=1.5cm,includeheadfoot,includehead,includefoot]{geometry}
\hypersetup{colorlinks,linkcolor=black}
\usepackage{fancyhdr}
\pagestyle{fancyplain}
\chead{}
\lhead{}
\rhead{}
\cfoot{}
\lfoot{\begin{footnotesize}alvaro.gonzalezsotillo@educa.madrid.org\end{footnotesize}}
\rfoot{\begin{footnotesize}\thepage / \pageref{LastPage}\end{footnotesize}}
\usepackage[skins]{tcolorbox}
\usepackage{multicol}
\usepackage{changepage} %ajdustwidth
\usepackage{fancybox}
\usepackage{attachfile2}
\lhead{Extraordinaria 2021 (es el \\lhead)}
\rhead{Administración y Gestión de Bases de Datos (es el \\rhead)}
\lhead{ED}
\rhead{Práctica metodologías ágiles}
\usepackage{svg}
\usepackage{letltxmacro}
\LetLtxMacro{\originalincludegraphics}{\includegraphics}
\renewcommand{\includegraphics}[2][]{\IfFileExists{#2.pdf}{\originalincludegraphics[#1]{#2.pdf}}{\originalincludegraphics[#1]{#2}}}
\LetLtxMacro{\originalincludesvg}{\includesvg}
\renewcommand{\includesvg}[2][]{\IfFileExists{#2.pdf}{\originalincludegraphics[#1]{#2.pdf}}{\originalincludegraphics[#1]{#2.svg.pdf}}}
\usepackage{comment}
\excludecomment{NOTES}
\date{}
\title{Metodologías ágiles}
\hypersetup{
 pdfauthor={Álvaro González Sotillo},
 pdftitle={Metodologías ágiles},
 pdfkeywords={},
 pdfsubject={},
 pdfcreator={Emacs 30.2 (Org mode 9.7.11)}, 
 pdflang={Spanish}}
\begin{document}

\maketitle
\setcounter{tocdepth}{1}
\tableofcontents

\captionsetup{font=scriptsize}

\textattachfile{/datos-luks/datos-alvaro/apuntes-clase/entornos-de-desarrollo/apuntes/1/ED-01-practica-metodologias-agiles.org}{} \vspace{-1em}


\setlength{\parindent}{0em}
\setlength{\parskip}{1em}

\newtcolorbox{Aviso}[1][Aviso]{
  enhanced,
  colback=gray!5!white,
  colframe=gray!75!black,fonttitle=\bfseries,
  colbacktitle=gray!85!black,
  attach boxed title to top left={yshift=-2mm,xshift=2mm},
  title=#1
}

\newtcolorbox{cuadrito}[1][Ignorado]{
  %drop shadow southeast,
  enhanced jigsaw,
  colback=white,
}


\newcommand{\StudentData}{
  \begin{cuadrito}[1\textwidth]
    \vspace{0.3cm}
    \large{
      \textbf{Apellidos:} \hrulefill \\
      \textbf{Nombre:} \hrulefill \\
      \textbf{Fecha:} \hrulefill \hspace{1cm} \textbf{Usuario:} \hrulefill \\
    }
    \vspace{-0.2cm}
  \end{cuadrito}
}

\newcommand{\StudentDataDniFirma}{
  \begin{cuadrito}[1\textwidth]
    \vspace{0.3cm}
    \large{
      \textbf{Apellidos:} \hrulefill \\
      \textbf{Nombre:} \hrulefill  \\
      \textbf{Dni:} \hrulefill \hspace{8cm} \\
      \\
      \textbf{Firma:} \hrulefill \hspace{8cm} \\
    }
    \vspace{-0.2cm}
  \end{cuadrito}
}
\section*{Objetivos de la práctica:}
\label{sec:org0000000}
En esta práctica se espera que el alumno alcance los siguientes objetivos:
\begin{itemize}
\item Se han identificado las fases de desarrollo de una aplicación informática.
\item Se ha evaluado la funcionalidad ofrecida por las herramientas utilizadas en el desarrollo de software.
\item Se han identificado las características y escenarios de uso de las metodologías ágiles de desarrollo de software.
\end{itemize}
\section{Elección de grupo y herramienta}
\label{sec:org0000003}
Se formarán grupos de, preferiblemente, tres alumnos. Cada grupo elegirá una herramienta asociada a las metodologías ágiles:
\begin{itemize}
\item Lean, Design sprint, Agile inception, XP, SCRUM, KANBAN, Test-Driven Development (TDD), Continuous Integration, Pair Programming\ldots{}
\end{itemize}

Formaliza \href{https://docs.google.com/spreadsheets/d/1-r-t4JHkw15FHfjQbuXN0CBH2zV-dWmBIPGqJIP8ZrQ/edit?gid=0\#gid=0}{tu grupo en esta hoja de cálculo}  
\section{En qué consiste el trabajo}
\label{sec:org0000006}
Después de haber elegido cada pareja/grupo una herramienta que utiliza metodologías ágiles, se procederá a investigar sobre ella para proceder a una exposición en clase sobre la misma, 

Podrás hacer uso de una presentación power point, canva, o similar. Este recurso de presentación se enviará a la tarea del aula virtual por todos los integrantes del equipo. En el caso de un recurso online, se enviará un documento con la URL.

El día del la presentación, el profesor abrirá el recurso entregado en el proyector, y cederá su sitio al grupo. Se dispondrá de 9 minutos para la exposición.
\section{Qué se valorará}
\label{sec:org0000009}

\begin{itemize}
\item Individual (40\%)
\begin{itemize}
\item (30\%) \textbf{Contenido}. Se nota un buen dominio de la parte que le ha tocado exponer, no comete errores, no duda.
\item (10\%) \textbf{Exposición}. Atrae la atención del público y mantiene el interés durante toda la exposición.
\end{itemize}

\item Grupal (60\%)
\begin{itemize}
\item (40\%) \textbf{Contenido}: El material preparado permite entender la herramienta elegida. Bien resumido, ordenado, claro, con los diagramas necesarios para ser entendido.
\item (10\%) \textbf{Agile}. Se establecerá la relación con el \emph{agile manifiesto}
\item (10\%) \textbf{Trabajo en equipo}. La exposición muestra planificación y trabajo en equipo reflejando que todos han colaborado de forma equitativa y organizada. Todos los miembros exponen por igual.
\end{itemize}
\end{itemize}
\section*{Resultados de aprendizaje}
\label{sec:org000000c}
Esta práctica contribuye a la calificación de los siguientes RA:
\begin{itemize}
\item RA 1
\end{itemize}
\end{document}
